\documentclass[11pt,a4paper]{article}
\usepackage{amsmath,amssymb,amsthm}
\usepackage[margin=2.5cm]{geometry}
\usepackage{hyperref}

\newtheorem{definition}{Definition}[section]
\newtheorem{theorem}[definition]{Theorem}
\newtheorem{proposition}[definition]{Proposition}
\newtheorem{lemma}[definition]{Lemma}
\newtheorem{corollary}[definition]{Corollary}
\newtheorem{conjecture}[definition]{Conjecture}
\newtheorem{remark}[definition]{Remark}

\title{Hyperseed Formalization:\\
  Big Tech's Big Bets on the Singularity}
\author{ProtomegaTron (automated formalization)}
\date{2026-09-17}

\begin{document}
\maketitle

\begin{abstract}
We formalize the intellectual structure of Ben Goertzel's Substack article
``Big Tech's Big Bets on the Singularity --- Will They Pay Off, or Blow Up?''
(2026-08-02) within the Hyperseed ontology. The article presents a detailed
economic scenario analysis of the AI industry's circular financing structure,
arguing that the over-bet is not a failure of prudence but the Nash
equilibrium of a US--China geopolitical prisoner's dilemma. Each structural
pattern maps onto Hyperseed structures: the prisoner's dilemma as cocycle
lock-in, the financial crash as base-space perturbation orthogonal to
the technology fiber, the decentralized network as fractional-Kelly fiber,
fire-sale hardware absorption as crisis-driven fiber elasticity, and the
commoditization paradox as a global section that doesn't project onto
the financial fiber.
\end{abstract}

\tableofcontents

%% =================================================================
\section{Circular financing and the bridge-or-kite question}
\label{sec:circular}
%% =================================================================

\begin{definition}[Bridge vs.\ kite --- claims 1--2]
\label{def:bridge}
Check-kiting and bridge financing are structurally identical: a circular
financing structure where each node looks funded by deposits from others.
The difference is determined only in retrospect:
\[
  \text{Bridge financing} \iff \text{deposit arrives before checks clear}
\]
\[
  \text{Kiting} \iff \text{deposit never arrives}
\]
The AI industry's circular deal structure --- Nvidia $\to$ hyperscalers
$\to$ cloud $\to$ startups $\to$ Nvidia --- has Ponzi-esque dynamical
characteristics. Whether it is bridge financing or kiting depends entirely
on whether the technology delivers revenue that breaks the circle.
\end{definition}

%% =================================================================
\section{The geopolitical prisoner's dilemma as cocycle lock-in}
\label{sec:pd}
%% =================================================================

\begin{proposition}[Aggregation and over-Kelly --- claims 4--7, 35]
\label{prop:aggregate}
The over-bet has two layers:
\begin{enumerate}
\item \textbf{Aggregation problem}: each firm sizes its bet rationally
  against its own balance sheet, but the deal web correlates everyone
  with everyone. The aggregate position is a collectively over-Kelly
  bet on a single coin flip, with no actor managing network-level
  exposure.
\item \textbf{Prisoner's dilemma}: the US and China perceive a
  winner-take-most race for the most strategically consequential
  technology in history. Prudent sizing is strategically dominated:
  \[
    (\text{size down}, \text{size down}) \succ (\text{all-in}, \text{all-in})
    \quad\text{but}\quad
    (\text{all-in}, \text{all-in}) = \text{Nash equilibrium}
  \]
  Contest theory predicts that evenly matched rivals dissipate nearly
  the entire prize through competitive expenditure.
\end{enumerate}

In Hyperseed terms, this is \emph{cocycle lock-in}: the race dynamics
lock both actors into a cocycle where defection (over-betting) is the
only stable transition function. Cooperation would require a different
cocycle (coordinated draw), but the levered financial structure is
itself priced against coordination --- even peace is a systemic risk
(claim 34) because removing the race logic collapses the justification
for the spending.
\end{proposition}

\begin{proposition}[Two disease architectures --- claims 8--10]
\label{prop:diseases}
The overinvestment manifests differently in the two systems:
\begin{itemize}
\item \textbf{American}: cardiac failure (margin calls, fire sales,
  acute crisis),
\item \textbf{Chinese}: metabolic failure (overcapacity, deflation,
  chronic absorption by state banks).
\end{itemize}
Same overinvestment, two pathologies. TSMC/Taiwan is the covariance
term connecting both --- the single physical asset on which both
national bets are jointly written.
\end{proposition}

%% =================================================================
\section{Financial crisis as base-fiber orthogonality}
\label{sec:crash}
%% =================================================================

\begin{theorem}[Modal and uninformative --- claims 11--13, 32, 39]
\label{thm:modal}
A financial crisis in 2027--2030 is the modal outcome under all three
technology timelines: fast (HLAGI-2029), moderate (HLAGI-2035), and
pessimist (plateau). Therefore:
\[
  P(\text{crash} \mid \text{fast}) \approx
  P(\text{crash} \mid \text{moderate}) \approx
  P(\text{crash} \mid \text{pessimist}) \approx \text{high}
\]
By Bayes' theorem, the crash carries almost no information about the
technology:
\[
  \frac{P(\text{fast} \mid \text{crash})}{P(\text{fast})} \approx 1
\]

In Hyperseed terms: the crash operates in the economic base space $B$,
which is approximately orthogonal to the technology fiber $F$. Financial
failure is a base-space perturbation that says nothing about the fiber's
structure. Leverage adjudicates timing; it says nothing about the terminus.

The Kurzweilian world and the Marcus world are observationally similar
for the next 24 months --- a fact worth stating loudly.
\end{theorem}

%% =================================================================
\section{Six scenarios: three timelines $\times$ two variants}
\label{sec:scenarios}
%% =================================================================

\begin{proposition}[Scenario 1: HLAGI-2029 without deAGI --- claims 14--18]
\label{prop:s1}
Even in the winning scenario, three problems:
\begin{enumerate}
\item \textbf{Commoditization}: if intelligence commoditizes (several
  labs at frontier, open models close behind), enormous consumer surplus
  coexists with poor returns on depreciating GPU fleets. The global
  section of the value sheaf exists but doesn't project onto the
  financial fiber (claim 40).
\item \textbf{Demand crisis}: severe wage displacement before
  redistribution machinery exists --- a demand-side collapse inside
  a supply-side miracle.
\item \textbf{Geopolitical instability}: two states approaching ASI at
  parity is classically unstable --- mobilization, expropriation by
  victory rather than bankruptcy by defeat.
\end{enumerate}
\end{proposition}

\begin{proposition}[Scenario 2: HLAGI-2035 without deAGI --- claims 19--21]
\label{prop:s2}
Revenue climbs but doesn't cross the obligation line until $\sim$2032.
The 2028--2030 refunding gate falls before the payoff arrives --- a
classic infrastructure crash where the technology was right but the
financing structure kills the pioneers (dot-com pattern). Too-strategic-to-fail
bailouts convert private loss to public debt.
\end{proposition}

\begin{proposition}[Scenario 3: Plateau without deAGI --- claims 22--23]
\label{prop:s3}
Revenue never crosses obligations. Permanent shortfall. Worst case for
the levered complex. Crash fuels blame politics and diversionary
nationalism.
\end{proposition}

%% =================================================================
\section{The decentralized variant: fractional-Kelly fiber}
\label{sec:deagi}
%% =================================================================

\begin{definition}[Fractional-Kelly architecture --- claims 30--31, 37]
\label{def:kelly}
The decentralized AGI network is intrinsically fractional-Kelly:
\begin{itemize}
\item No debt service (no leveraged capex to refinance),
\item No frontier-hardware treadmill (runs on commodity and distressed
  hardware),
\item Elastic capacity (pay-as-you-go token flows),
\item Modest continuous stakes with no ruin branch.
\end{itemize}
In fiber terms: the risk is distributed across the fiber rather than
concentrated in a single highly-leveraged node. The fiber structure
$E_{\text{deAGI}}$ has no ruin singularity:
\[
  P(\text{ruin} \mid E_{\text{deAGI}}) \approx 0
  \quad\text{vs.}\quad
  P(\text{ruin} \mid E_{\text{levered}}) \gg 0
\]
\end{definition}

\begin{proposition}[Crisis absorption --- claims 26--29, 38]
\label{prop:absorption}
In Scenario 2 + deAGI, the crash feeds the network:
\begin{enumerate}
\item Fire-sale hardware from bankrupt hyperscalers jumps network
  capacity --- the fiber stretches when the base space contracts
  (crisis absorption via fiber elasticity).
\item Too-strategic-to-fail bailout logic weakens when the frontier
  is demonstrably elsewhere.
\item Non-aligned countries prefer a neutral network (no alignment
  strings) --- third-market default.
\item Capex shifting to pay-as-you-go flows stops AI from crowding
  the Treasury market --- macro dividend via disintermediation.
\end{enumerate}
\end{proposition}

\begin{theorem}[Fractional Kelly wins Scenario 3 --- claim 31]
\label{thm:kelly}
In the pessimist scenario where all AI plateaus, the decentralized
network wins not by being right about AGI but by being the only
structure that didn't bet its existence on it. In a commoditized
narrow-AI world of thin margins, the minimum-cost producer is the
one with no debt service and elastic capacity --- the network.

This is the Kelly criterion applied to institutional design:
Full Kelly maximizes expected log-growth but has nonzero ruin
probability; fractional Kelly sacrifices peak growth rate for
survival. The states went Full Kelly; the network went fractional.
\end{theorem}

%% =================================================================
\section{Peace as systemic risk}
\label{sec:peace}
%% =================================================================

\begin{proposition}[Peace destabilizes --- claims 33--34, 41]
\label{prop:peace}
Even peace is a systemic risk in the current architecture. A negotiated
draw removes the race logic that justifies the spending, collapsing the
revenue expectations that the levered structure is priced on. The
financial structure is stable only under adversarial cocycle conditions:
\[
  \text{stability}(E_{\text{levered}}) \iff
  \text{adversarial cocycle persists}
\]
This creates a perverse dynamic where the financial system has a
structural interest in continued rivalry, independent of the wishes
of the actors within it.
\end{proposition}

%% =================================================================
\section{The commoditization paradox}
\label{sec:commoditization}
%% =================================================================

\begin{proposition}[Consumer surplus without producer surplus --- claim 40]
\label{prop:commodity}
The technology succeeds (consumer surplus enormous) while the investors
fail (producer surplus near zero). In Hyperseed terms:
\[
  \exists\; \text{global section } s : B \to E_{\text{value}}
  \quad\text{but}\quad
  \pi_{\text{financial}}(s) \approx 0
\]
The global section of the value sheaf exists --- society captures the
value of AGI --- but it doesn't project onto the financial fiber ---
investors see no return. This is the airline paradox at civilizational
scale: transformative technology with century-long aggregate industry losses.
\end{proposition}

%% =================================================================
\section{Cross-references}
%% =================================================================

\begin{remark}[Connections to other formalizations]
\begin{itemize}
\item \textbf{Geopolitics} (2026-08-24): extends the same 54-cell grid
  analysis with detailed political branching. This article provides the
  economic detail; the Geopolitics article provides the political detail.
  Financial crisis as base-space perturbation with political holonomy
  is the shared structure.
\item \textbf{Zuck's Future} (2026-08-10): access vs.\ ownership and
  correct fiber bundle architecture. The decentralized alternative here
  is the same architecture Zuck claims to offer but doesn't.
\item \textbf{How Omega Lost Its Claw} (2026-09-11): distributed
  development as fractional-Kelly. The engineering architecture of
  decentralized AGI is described there; the economic resilience is
  described here.
\item \textbf{Seven Flavors} (2026-07-01): arms race as cross-cutting
  attractor. Here the arms race dynamics are given their economic
  mechanism (prisoner's dilemma over bet sizing).
\item \textbf{Sanders Ban} (2026-09-05): state capture of AI development.
  The too-strategic-to-fail bailout pathway here is the economic version
  of the Sanders ban pathway --- state capture via financial rescue
  rather than legislative prohibition.
\item \textbf{Provenance Layer} (2026-08-25): decentralized resolution
  as distributed fiber bundle. Same anti-monoculture principle applied
  to economic infrastructure.
\end{itemize}
\end{remark}

\begin{conjecture}[Crisis-to-network conversion rate]
Conjecture: the rate at which a financial crisis feeds the decentralized
network is proportional to the price gap between distressed hardware and
network marginal cost. When the gap is large (deep crisis), hardware
flows rapidly into the network; when small (mild correction), the
conversion stalls. The network's crisis absorption capacity is therefore
a function of its current utilization rate and the severity of the crisis
--- and the network should be designed with excess absorption capacity
in anticipation of the modal crisis.
\end{conjecture}

\end{document}
